
Diante de um cenário nacional de escassez e estiagens hídricas, a mitigação de desperdícios torna-se uma tarefa de extrema urgência \cite{tratabrasil2025}.

A digitalização da medição de consumo hídrico apresenta-se como uma ferramenta estratégica na contenção de desperdícios e compreensão de padrões de uso. 
Entretanto, a modernização da atual infraestrutura enfrenta desafios técnicos e financeiros, especialmente no que tange à manutenção, confiabilidade e custo no médio/longo prazo sob os dispositivos empregados em tal tarefa.

Neste contexto, este trabalho tem como objetivo comparar duas formas de implementação de sistemas telemétricos baseados em IOT, mais especificamente, com computação centralizada (Cloud Computing) e com computação em borda (Edge Computing), com o objetivo de definir qual a melhor abordagem na resolução do problema.